\documentclass[12pt,a4paper]{amsart}

\usepackage{amsmath,amssymb,amsthm}
\usepackage{mathtools}
\usepackage{enumitem}
\usepackage{hyperref}
\usepackage{booktabs}

%%%─ Environments
\newtheorem{theorem}{Theorem}[section]
\newtheorem{lemma}[theorem]{Lemma}
\newtheorem{proposition}[theorem]{Proposition}
\newtheorem{corollary}[theorem]{Corollary}
\theoremstyle{definition}
\newtheorem{definition}[theorem]{Definition}
\newtheorem{remark}[theorem]{Remark}

%%%─ Macros (reuse from XVIII)
\newcommand{\cA}{\mathcal{A}}
\newcommand{\R}{\mathbb{R}}
\newcommand{\cK}{K_{B_c}}
\newcommand{\cKA}{K_{\cA}}
\newcommand{\Ai}{\mathrm{Ai}}
\newcommand{\PK}{P_K}
\newcommand{\PcK}{\widetilde{P}_K^{(c)}}
\newcommand{\Uc}{\widetilde{U}_c}
\newcommand{\Lc}{\mathcal{L}_c^{\mathrm{resc}}}
\newcommand{\xit}{\tilde{\xi}}
\newcommand{\normh}[2]{\|#1\|_{#2}}
\newcommand{\qa}{\mathfrak{a}}
\newcommand{\SK}{S_K}
\newcommand{\HS}{\mathrm{HS}}
%%% Paper XIX specific
\newcommand{\Cprime}{C'(M,K)}
\newcommand{\ConeM}{C_1(M,K)}
\newcommand{\Clin}{C_{\mathrm{lin}}(M,K)}
\newcommand{\Cquad}{C_{\mathrm{quad}}(M,K)}
\newcommand{\Ctwo}{C_2(\alpha)}
\newcommand{\Kopt}{K_{\mathrm{opt}}}
\newcommand{\AK}{A_K}
\newcommand{\BK}{B_K}
\newcommand{\Cinfty}{C_\infty(M,K)}

\subjclass[2020]{47B34, 45C05, 34E20, 47G10, 60B20}
\keywords{Prolate spheroidal wave functions, Airy kernel,
  two-scale convergence, quantitative rate, Langer error,
  spectral tail, polynomial truncation cost, BR3}

\begin{document}

\title[Paper XIX: Quantitative Rate for BR3]{%
  Paper~XIX: Quantitative Convergence Rate\\[0.3em]
  \large for Airy Kernel Universality (BR3)\\[0.2em]
  \normalsize\textsc{Freeze Version} --- Two-Scale Structure Closed}

\author{\textit{prolate-gram-coercivity programme}}
\date{May 2026}

\begin{abstract}
Building on Paper~XVIII (FREEZE~\texttt{d7609d2}),
which established $\normh{\cK - \cKA}{L^2(w\otimes w)}\to 0$ (BR3),
this paper supplies explicit convergence rates
via a \emph{two-scale decomposition}
\begin{equation}\label{eq:decomp_abstract}
  \SK - \cKA = \AK - \BK,
\end{equation}
where $\AK$ (Langer regime, rank $\leq 3K$) and
$\BK$ (spectral tail) are controlled by entirely disjoint mechanisms
with vertex-disjoint dependency subgraphs.

The main structural results are:
\begin{itemize}[nosep]
  \item \textbf{Lemma~\ref{lem:langer_rate}} (Langer scale,
    fixed $K$, $c\to\infty$):
    \[
      \normh{\AK}{L^2([0,M]^2)}
      \leq \Clin\cdot c^{-1/3} + \Cquad\cdot c^{-2/3},
    \]
    where $\Clin = O(MK^{5/3})$ (linear interaction: one exact
    eigenfunction, one Langer error) and
    $\Cquad = O(M^2K^{7/3})$ (quadratic: two Langer errors).
    The linear term dominates for $c\geq 1$;
    the quadratic term is subleading by an additional $c^{-1/3}$ factor.
  \item \textbf{Lemma~\ref{lem:tail_rate}} (spectral scale,
    $K\to\infty$, uniform in $c$):
    $\normh{\BK}{L^2([0,M]^2)}\leq C_2(\alpha)^{1/2}\cdot e^{-\alpha K/2}$.
\end{itemize}
These combine to a two-scale bound
(Theorem~\ref{thm:twoscale}) with no joint uniformity in
$(c,K)$ assumed or needed.
The coupling $\Kopt(c) = \lfloor(\log c)/(3\alpha)\rfloor$
yields the rate $O(c^{-1/3}(\log c)^{5/3})$
(Corollary~\ref{cor:lograte}), where the quadratic remainder
$O(c^{-2/3}(\log c)^{7/3})$ is strictly subleading.
\end{abstract}

\maketitle
\tableofcontents
\bigskip\hrule\bigskip

%% ============================================================
\section{Setup and Inherited Notation}
\label{sec:setup}
%% ============================================================

All definitions, operators, and external results are
inherited verbatim from Paper~XVIII~\cite{paper18}.
We introduce only the central decomposition and its constants.

\begin{definition}[Two-scale error decomposition]
\label{def:decomp}
For fixed $M < \infty$, $K \geq 1$, and $c \geq c_0(M,K)$, write
$e_j := \tilde\psi_j - u_j \in C^0([0,M])$
(XVIII Lemma~\texttt{lem:H2}). Define:
\begin{align}
  \AK(\xit,\xit') &:=
    \sum_{j=1}^K\bigl(
      e_j(\xit)\,u_j(\xit') + u_j(\xit)\,e_j(\xit') + e_j(\xit)\,e_j(\xit')
    \bigr),
    \label{eq:AK} \\
  \BK(\xit,\xit') &:= \sum_{j>K} u_j(\xit)\,u_j(\xit').
    \label{eq:BK}
\end{align}
Then algebraically $\SK - \cKA = \AK - \BK$.
\end{definition}

\begin{definition}[Langer-regime constants]
\label{def:langerconstants}
For fixed $M,K$, let:
\begin{align}
  \Cprime &:= \text{constant from XVIII~\texttt{lem:H2}}:
    \quad\normh{e_j}{C^0([0,M])} \leq \Cprime\,c^{-1/3},
    \quad j\leq K; \label{eq:Cprime}\\
  \Cinfty &:= \sup_{j\leq K}\normh{u_j}{C^0([0,M])}
    = O(1) \quad\text{(XVIII~\texttt{lem:linf})}; \label{eq:Cinfty}\\
  \Clin   &:= 2KM\,\Cprime\,\Cinfty
    = O(MK^{5/3}); \label{eq:Clin} \\
  \Cquad  &:= KM\,(\Cprime)^2
    = O(M^2K^{7/3}). \label{eq:Cquad}
\end{align}
By XVIII~\texttt{rem:langerkdep}: $\Cprime = O(MK^{2/3})$,
so $\Clin = O(MK^{5/3})$ and $\Cquad = O(M^2K^{7/3})$.
Both are independent of $c$.
\end{definition}

\begin{remark}[Erratum to original $C_1(M,K)$ definition]
\label{rem:C1scale}
The original working version of this paper wrote
$C_1(M,K) := KM\Cprime(2\Cinfty+\Cprime) = O(M^2K^{5/3})$
and cited $O(K^{5/3})$ as the truncation cost.
This is \emph{asymptotically correct at the level of the rate}
$O(c^{-1/3}(\log c)^{5/3})$,
but the internal constant is imprecise:
the quadratic term $KM(\Cprime)^2 = O(K^{7/3})$
is not $O(K^{5/3})$.

The refined bookkeeping (Definition~\ref{def:langerconstants} above)
splits the contribution of $\AK$ into two terms:
\[
  \normh{\AK}{} \leq \Clin\cdot c^{-1/3} + \Cquad\cdot c^{-2/3},
\]
where $\Clin = O(MK^{5/3})$ from linear interactions
($u_j \otimes e_j + e_j \otimes u_j$)
and $\Cquad = O(M^2K^{7/3})$ from the quadratic interaction
($e_j \otimes e_j$, carrying an extra $c^{-1/3}$ from the
second Langer error factor).
At $K = \Kopt \sim (\log c)/(3\alpha)$ the quadratic term is
$O(c^{-2/3}(\log c)^{7/3}) = o(c^{-1/3}(\log c)^{5/3})$,
confirming that the rate of Corollary~\ref{cor:lograte}
is correct. Only the single-constant formulation was imprecise.
\end{remark}

\begin{definition}[Slepian tail constant]
\label{def:C2}
Let $\alpha = \min(\beta_1,\beta_2) > 0$ from XVIII~\texttt{lem:slepianunif}.
Set $\Ctwo := (1-e^{-\alpha})^{-1}$.
\end{definition}

%% ============================================================
\section{Dependency Graph}
\label{sec:depgraph}
%% ============================================================

\begin{center}
\renewcommand{\arraystretch}{1.2}
\begin{tabular}{l l l l}
\toprule
\textbf{Node} & \textbf{Introduces} & \textbf{Internal deps}
  & \textbf{From $\mathcal{E}$} \\
\midrule
$d_0$: Def.~\ref{def:decomp}
  & $\AK$, $\BK$, $e_j$ & ---
  & XVIII lem:H2 \\
$d_1$: Def.~\ref{def:langerconstants}
  & $\Clin$, $\Cquad$, $\Cprime$, $\Cinfty$ & $d_0$
  & XVIII lem:H2, lem:linf, \textbf{rem:langerkdep} \\
$d_2$: Def.~\ref{def:C2}
  & $\Ctwo$, $\alpha$ & ---
  & XVIII lem:slepianunif \\
\hline
$\ell_1$: Lem.~\ref{lem:langer_rate}
  & $\normh{\AK}{}\leq\Clin c^{-1/3}+\Cquad c^{-2/3}$ & $d_0,d_1$
  & XVIII lem:H2, lem:linf \\
$\ell_2$: Lem.~\ref{lem:tail_rate}
  & $\normh{\BK}{}\leq\Ctwo^{1/2}e^{-\alpha K/2}$ & $d_0,d_2$
  & XVIII lem:besseltail, lem:slepianunif \\
\hline
$T$: Thm.~\ref{thm:twoscale}
  & two-scale bound & $\ell_1,\ell_2$ & --- \\
$C$: Cor.~\ref{cor:lograte}
  & $O(c^{-1/3}(\log c)^{5/3})$ & $T,d_1$ & --- \\
\bottomrule
\end{tabular}
\end{center}

%% ============================================================
\section{Main Lemmas}
\label{sec:lemmas}
%% ============================================================

\subsection{Level 1: Langer regime (fixed $K$, $c\to\infty$)}
\label{ssec:langer}

\begin{lemma}[Langer-regime rate, refined]
\label{lem:langer_rate}
Fix $M < \infty$ and $K \geq 1$.
For all $c \geq c_0(M,K)$:
\begin{equation}\label{eq:term1}
  \normh{\AK}{L^2([0,M]^2)}
  \leq \Clin\cdot c^{-1/3} + \Cquad\cdot c^{-2/3},
\end{equation}
where $\Clin = O(MK^{5/3})$ and $\Cquad = O(M^2K^{7/3})$
are independent of $c$.
\end{lemma}

\begin{proof}
Fix $j\leq K$. Using
$\normh{f\otimes g}{L^2([0,M]^2)} \leq M\normh{f}{C^0}\normh{g}{C^0}$
together with~\eqref{eq:Cprime}--\eqref{eq:Cinfty}:
\begin{align}
  \normh{e_j\otimes u_j}{L^2([0,M]^2)}
    &\leq M\Cprime\Cinfty\,c^{-1/3}, \label{eq:c1}\\
  \normh{u_j\otimes e_j}{L^2([0,M]^2)}
    &\leq M\Cinfty\Cprime\,c^{-1/3}, \label{eq:c2}\\
  \normh{e_j\otimes e_j}{L^2([0,M]^2)}
    &\leq M(\Cprime)^2 c^{-2/3}. \label{eq:c3}
\end{align}
Note that~\eqref{eq:c3} carries $c^{-2/3}$, \emph{not} $c^{-1/3}$:
the quadratic interaction contributes at a strictly lower $c$-scale.
Triangle inequality per $j$ and summing over $j=1,\ldots,K$:
\[
  \normh{\AK}{L^2([0,M]^2)}
  \leq \underbrace{2KM\Cprime\Cinfty}_{=\,\Clin}\cdot c^{-1/3}
  \;+\; \underbrace{KM(\Cprime)^2}_{=\,\Cquad}\cdot c^{-2/3}.
\]
\end{proof}

\begin{remark}[Two internal scales inside $\AK$]
\label{rem:twointernal}
Lemma~\ref{lem:langer_rate} exhibits a two-scale structure
\emph{within} the Langer term $\AK$ itself:
\begin{itemize}[nosep]
  \item \emph{Linear interaction} ($u_j\otimes e_j + e_j\otimes u_j$):
    one exact eigenfunction ($C_\infty$) times one Langer error ($\Cprime c^{-1/3}$);
    contributes $O(K^{5/3})\cdot c^{-1/3}$.
  \item \emph{Quadratic interaction} ($e_j\otimes e_j$):
    two Langer errors; contributes $O(K^{7/3})\cdot c^{-2/3}$.
\end{itemize}
At $K = \Kopt \sim (\log c)/(3\alpha)$:
$O(K^{7/3})c^{-2/3} = O((\log c)^{7/3}c^{-2/3}) = o(c^{-1/3}(\log c)^{5/3})$,
so the linear interaction dominates and the asymptotic rate is $O(c^{-1/3}(\log c)^{5/3})$.
\end{remark}

\begin{remark}[Polynomial truncation cost is geometrically inevitable]
\label{rem:polycost}
The growth $\Clin = O(MK^{5/3})$ reflects Airy zero spacing
$a_K\sim K^{2/3}$ propagating through the coercivity spine
(XVIII~\texttt{rem:langerkdep}).
This polynomial cost does not affect the two-scale limit
structure at fixed $K$.
\end{remark}

\subsection{Level 2: Spectral tail (uniform in $c$, $K\to\infty$)}
\label{ssec:tail}

\begin{lemma}[Spectral tail rate]
\label{lem:tail_rate}
For $K \geq K_0$ and all $c \geq c_1$:
\begin{equation}\label{eq:term2}
  \normh{\BK}{L^2([0,M]^2)}
  \leq \Ctwo^{1/2}\cdot e^{-\alpha K/2}.
\end{equation}
The bound is uniform in $c$.
\end{lemma}

\begin{proof}
By Cauchy--Schwarz:
$\normh{\BK}{L^2([0,M]^2)}^2
\leq\bigl(\sum_{j>K}\normh{u_j}{L^2([0,M])}^2\bigr)^2$.
Slepian--Landau (XVIII~\texttt{lem:slepianunif}):
$\normh{u_j}{L^2([0,M])}^2 \leq e^{-\alpha j}$.
Hence $\sum_{j>K}\normh{u_j}{L^2}^2 \leq \Ctwo e^{-\alpha K}$;
taking square roots gives~\eqref{eq:term2}.
\end{proof}

%% ============================================================
\section{Main Theorem and Corollary}
\label{sec:main}
%% ============================================================

\begin{theorem}[Two-scale convergence rate for BR3]
\label{thm:twoscale}
For fixed $M < \infty$, $K \geq K_0$, and
$c \geq \max(c_0(M,K),\,c_1)$:
\begin{equation}\label{eq:twoscale}
  \normh{\cK - \cKA}{L^2([0,M]^2)}
  \leq
  \Clin\cdot c^{-1/3}
  + \Cquad\cdot c^{-2/3}
  + \Ctwo^{1/2}\cdot e^{-\alpha K/2}.
\end{equation}
\end{theorem}

\begin{proof}
Triangle inequality on $\AK - \BK$ + Lemmas~\ref{lem:langer_rate}
and~\ref{lem:tail_rate}.
\end{proof}

\begin{corollary}[Rate $O(c^{-1/3}(\log c)^{5/3})$]
\label{cor:lograte}
Let $\Kopt(c) := \lfloor(\log c)/(3\alpha)\rfloor$.
For $c$ large enough that $\Kopt\geq K_0$:
\begin{equation}\label{eq:lograte}
  \normh{\cK - \cKA}{L^2([0,M]^2)}\Big|_{K=\Kopt}
  = O\bigl(c^{-1/3}(\log c)^{5/3}\bigr).
\end{equation}
\end{corollary}

\begin{proof}
At $K = \Kopt \sim (\log c)/(3\alpha)$:
\begin{itemize}[nosep]
  \item \emph{Linear term:}
    $\Clin c^{-1/3} = O(K^{5/3})c^{-1/3} = O(c^{-1/3}(\log c)^{5/3})$.
  \item \emph{Quadratic term:}
    $\Cquad c^{-2/3} = O(K^{7/3})c^{-2/3}
    = O(c^{-2/3}(\log c)^{7/3}) = o(c^{-1/3}(\log c)^{5/3})$.
  \item \emph{Tail:}
    $e^{-\alpha\Kopt/2} \leq c^{-1/6} = o(c^{-1/3}(\log c)^{5/3})$.
\end{itemize}
The linear term dominates; \eqref{eq:lograte} follows.
\end{proof}

\begin{remark}[Interpretation of the log-loss]
\label{rem:logloss}
The factor $(\log c)^{5/3}$ traces Airy zero geometry:
$a_K\sim K^{2/3}$ drives $\Cprime = O(K^{2/3})$ and
$\Clin = O(K^{5/3})$ through the coercivity spine.
The quadratic interaction $\Cquad = O(K^{7/3})$ is suppressed
by an additional $c^{-1/3}$ and becomes subleading at $K = \Kopt$.
The log-exponent $5/3$ is the spectral-geometric invariant
of the linear interaction; this is made precise in Paper~XX.
\end{remark}

%% ============================================================
\section{Freeze Checklist}
\label{sec:checklist}
%% ============================================================

\begin{center}
\renewcommand{\arraystretch}{1.3}
\begin{tabular}{p{8.0cm} p{1.0cm} p{2.8cm}}
\toprule
\textbf{Claim} & \textbf{Status} & \textbf{Source} \\
\midrule
Algebraic decomposition $\SK-\cKA=\AK-\BK$
  & \checkmark & Def.~\ref{def:decomp} \\
Two-scale structural separation ($c\to\infty$ vs.\ $K\to\infty$)
  & \checkmark & \S\ref{sec:setup} \\
Lem.~\ref{lem:langer_rate}:
  $\normh{\AK}{}\leq\Clin c^{-1/3}+\Cquad c^{-2/3}$ (lin/quad split)
  & \checkmark & \S\ref{ssec:langer} \\
$\Clin=O(MK^{5/3})$, $\Cquad=O(M^2K^{7/3})$, both $c$-independent
  & \checkmark & Def.~\ref{def:langerconstants} \\
Quadratic term $\Cquad c^{-2/3}$ subleading at $K=\Kopt$
  & \checkmark & Rem.~\ref{rem:twointernal} \\
Lem.~\ref{lem:tail_rate}:
  $\normh{\BK}{}\leq\Ctwo^{1/2}e^{-\alpha K/2}$ (unif.\ in $c$)
  & \checkmark & \S\ref{ssec:tail} \\
Cor.~\ref{cor:lograte}: $O(c^{-1/3}(\log c)^{5/3})$ rigorous,
  linear term dominant
  & \checkmark & \S\ref{sec:main} \\
Erratum to original $C_1 = O(K^{5/3})$ bookkeeping
  & \checkmark & Rem.~\ref{rem:C1scale} \\
No joint $(c,K)$ uniformity assumed
  & \checkmark & \S\ref{sec:main} \\
\bottomrule
\end{tabular}
\end{center}

\begin{thebibliography}{99}
\bibitem{paper17} \textit{Paper~XVII}, 2026.
\bibitem{paper18} \textit{Paper~XVIII: Airy Kernel Universality},
  FREEZE commit \texttt{d7609d2}, 2026.
\bibitem{Kato1966} T.~Kato,
  \textit{Perturbation Theory for Linear Operators}, Springer, 1966.
\bibitem{ReedSimon2} M.~Reed, B.~Simon,
  \textit{Methods of Modern Mathematical Physics, Vol.~II},
  Academic Press, 1975.
\bibitem{Olver1997} F.~Olver,
  \textit{Asymptotics and Special Functions}, AK Peters, 1997.
\bibitem{TracyWidom1994} C.~Tracy, H.~Widom,
  Level-spacing distributions and the Airy kernel,
  \textit{Comm.\ Math.\ Phys.}\ \textbf{159} (1994), 151--174.
\bibitem{Slepian1978} D.~Slepian,
  Prolate spheroidal wave functions, Fourier analysis, and
  uncertainty~V,
  \textit{Bell Syst.\ Tech.\ J.}\ \textbf{57} (1978), 1371--1430.
\bibitem{Osipov2013} A.~Osipov, V.~Rokhlin, H.~Xiao,
  \textit{Prolate Spheroidal Wave Functions of Order Zero},
  Springer, 2013.
\end{thebibliography}

\end{document}
